% Layout extracted from Kupzyk report at ca1d062; retain a descriptive title.
\documentclass[11pt]{article}
\usepackage[letterpaper,margin=0.82in]{geometry}
\usepackage{booktabs,longtable,array,graphicx,amsmath,xcolor,hyperref,xurl}
\usepackage[authoryear,round]{natbib}
\hypersetup{colorlinks=true,urlcolor=blue,linkcolor=blue,citecolor=blue}
\renewcommand{\arraystretch}{1.25}
\setlength{\emergencystretch}{1em}
\setlength{\parskip}{5pt}\setlength{\parindent}{0pt}
% Use scalable fonts; inspect pdffonts output after building. Avoid adding T1
% fontenc without an installed scalable font family. Add title/date in manuscript.


\urlstyle{tt}
\newcolumntype{L}[1]{>{\raggedright\arraybackslash}p{#1}}
\newcommand{\disagreement}[1]{\par\textit{Disagreement record}\enspace\nolinkurl{#1}.}
\title{Simplifying assumptions in power analysis}
\author{}
\date{}
\makeatletter
\renewcommand{\@maketitle}{%
  \begin{center}
    {\LARGE\@title\par}
    \smallskip
    {\large A replication of Kupzyk's longitudinal cluster trial study\par}
  \end{center}}
\makeatother
\begin{document}
\maketitle

\begin{abstract}
We partly replicated the planning findings of \citet{kupzyk2011}, who reported that treatment imbalance and attrition cost more power than unequal cluster sizes, that losses were greatest at the lowest intraclass correlation, that up to eight extra clusters were needed, and that empirical power was usually lower than power calculated at the generating covariance. Exact replay requires the original software state, simulation streams, and dropout masks. The mean loss for 70/30 allocation was 6.409 versus 6.443 percentage points (published versus ours), and 15\% interval dropout cost 4.391 versus 4.472. The maximum extra count was eight versus seven, and empirical power was lower in 41 versus 40 of 54 cells. All 54 initial empirical powers were inside the adjusted comparison bands. Three sample-size crossings moved by one cluster, and the extreme-setting Optimal Design rejection rates differed.
\end{abstract}
\noindent\begin{minipage}{\linewidth}
\footnotesize
\setlength{\tabcolsep}{2.5pt}
\renewcommand{\arraystretch}{1.12}
\refstepcounter{table}\label{tab:headline}\textbf{Table \thetable. Headline comparisons.} Published findings and contrasts come from Kupzyk's Tables 1, 4 and 5 and pp. 106--107. Our values come from the R reconstruction.\par
\begin{tabular}{@{}L{79pt}L{75pt}L{75pt}r r L{48pt}L{48pt}l@{}}
\toprule
Contrast & Published & Ours & \multicolumn{2}{c}{Signed effect} & Effect magnitude & Direction & Verdict \\
 & & & Source & Ours & & & \\
\midrule
5\% dropout & 1.3 pp mean loss & 1.27 pp mean loss & 1.327 & 1.275 & smaller & same direction & match \\
15\% dropout & 4.4 pp mean loss & 4.47 pp mean loss & 4.391 & 4.472 & larger & same direction & match \\
70/30 allocation & 6.4 pp mean loss & 6.44 pp mean loss & 6.409 & 6.443 & larger & same direction & match \\
Normal sizes & 0.45 pp mean loss & 0.44 pp mean loss & 0.454 & 0.438 & smaller & same direction & match \\
Skewed sizes & 0.85 pp mean loss & 0.78 pp mean loss & 0.851 & 0.778 & smaller & same direction & match \\
ICC .05 vs .15 at 15\% dropout & All peak at .05. Skew rises .61 to .67 pp despite reported decline. & All peak at .05. Skew loss rises .56 to .62 pp. & 3.037 & 3.062 & larger & same direction & partial \\
Maximum extra clusters & Up to 8 extra clusters & Up to 7 extra clusters & 8 & 7 & smaller & same direction & partial \\
Count below minus 27 & 41 of 54 below & 40 of 54 below & 14 & 13 & smaller & same direction & match \\
\bottomrule
\end{tabular}
\par A positive effect denotes power lost versus the matched reference (percentage points, pp), greater dropout loss at ICC .05 than .15 (pp), extra clusters versus the between--within reference, or a count above half the cells, respectively. The rounding tolerances are 0.001 pp for the five main losses, 0.002 pp for the ICC contrast and 0.5 for integer counts. These point estimates do not establish equivalence. Verdicts address the scientific patterns. Monte Carlo, mask and reading uncertainty are assessed separately.
\end{minipage}

\section{Introduction}\label{sec:intro}
We partly replicated Kupzyk's headline findings. The estimates cover new samples from the reported design. Exact recovery of the original fits and realized dropout masks is unavailable. In the final planning grid, 70/30 allocation cost 6.409 versus 6.443 percentage points (published versus ours), and 5\% and 15\% interval attrition cost 1.327 versus 1.275 and 4.391 versus 4.472. Unequal sizes had smaller mean costs, and all five mean losses peaked at an intraclass correlation (ICC) of .05. Skewed sizes had a higher loss at .15 than at .10 in both analyses. We matched 51 of 54 cluster thresholds exactly. The maximum increase was eight versus seven clusters, and empirical power fell below exemplary power in 41 versus 40 of 54 cells.

The planning question is how much power a trial loses when convenient sample-size assumptions fail. \citet[pp.~1--6, 64--67]{kupzyk2011} compared three methods for a treatment-by-time interaction in longitudinal cluster trials. Formulaic power applies an analytic expression to a simplified design. Exemplary power uses the covariance and layout of a specified design. Empirical power is the fraction of repeated fitted trials that reject the null. These quantities are especially informative together when dropout changes the layout or finite-sample tests use different reference degrees of freedom.

\section{Power losses and sample-size consequences}\label{sec:effects}
Allocation imbalance had the largest mean power cost in both final exemplary grids, followed by high attrition. In those grids, the mean losses for normal and skewed cluster sizes were .454 and .851 pp in the published data and .438 and .778 pp in the reconstruction. Each contrast holds the other factors fixed. This gives 27 matched pairs for allocation and 18 for each attrition or size contrast. These are averages over the experimental settings. For example, one normal-size pair changes sign between the published and reconstructed final grids, and five do so in the initial empirical grid. In initial empirical cells 16 and 17 at ICC .05, the equal-minus-normal size contrast was $-.4183$ pp in the source and $-.11328$ pp in R, a small power gain for normal sizes in that pair. The small mean size effects describe the tested size distributions. The ranking applies to these four-occasion Gaussian designs and the specified dropout mechanism. See \citet[Tables 1--2 and 4]{kupzyk2011} and \href{https://github.com/alectionary/vibe-math/blob/master/formalization/kupzyk-power-assumptions/results/full/paired-contrasts.csv}{\nolinkurl{results/full/paired-contrasts.csv}}.

The published hypothesis placed attrition first, while the published results and the reconstruction place allocation first. The final-grid comparison is 6.409 versus 4.391 pp in the source and 6.443 versus 4.472 pp in R for allocation and 15\% attrition, respectively. The mean ordering holds when each final exemplary contrast changes just one factor, averaging over pairs matched on the other factors. The evidence explains the discrepancy as a result of the specific changes compared in this design (\href{https://github.com/alectionary/vibe-math/blob/master/formalization/kupzyk-power-assumptions/results/full/comparison-summary.json}{\nolinkurl{results/full/comparison-summary.json}}).
\label{dis:predicted-order}\disagreement{predicted-order}

The lowest ICC had the greatest loss for each violation, with one exception to a strictly declining trend across ICC. The skewed-size loss rises from .6098 to .6673 pp between ICC .10 and .15 in the source, and from .5614 to .6248 pp in R. For 15\% interval dropout, the loss at .05 exceeds the loss at .15 by 3.037 pp in the source and 3.062 pp in R. The source ordering at the lowest ICC is supported, and the monotonic generalization in the discussion fails under recomputation from Table 1 (\href{https://github.com/alectionary/vibe-math/blob/master/formalization/kupzyk-power-assumptions/results/source-identities.json}{\nolinkurl{results/source-identities.json}}).
\label{dis:skew-icc-monotonicity}\disagreement{skew-icc-monotonicity}

\subsection{How many clusters restore 80\% power?}\label{sec:thresholds}
The first threshold for mean exemplary power matched the published threshold in 51 of 54 settings. Three earlier crossings occurred just above .80. At ICC .05 with 5\% interval dropout, 70/30 allocation and normal sizes, 16 clusters replace the published 17, with mean power .80090. At ICC .05 with 15\% dropout, balanced allocation and equal sizes, 15 replace 16, with mean power .80086. At ICC .15 with 15\% dropout, 70/30 allocation and skewed sizes, 32 replace 33, with mean power .80190. At those three earlier counts, .360, .166 and .188 of masks fell below .80, respectively. Published threshold counts are compatible with a realized mask that crosses one cluster later. The largest extra count relative to the between--within reference was seven in R versus eight in Table 5. Neighboring counts and independent covariance checks support the three earlier mean crossings (\href{https://github.com/alectionary/vibe-math/blob/master/formalization/kupzyk-power-assumptions/results/full/threshold-comparison.csv}{\nolinkurl{results/full/threshold-comparison.csv}} and \href{https://github.com/alectionary/vibe-math/blob/master/formalization/kupzyk-power-assumptions/results/full/threshold-scan.csv}{\nolinkurl{results/full/threshold-scan.csv}}).
\label{dis:threshold-mask}\disagreement{threshold-mask}

\subsection{Oversampling for dropout}\label{sec:oversample}
Adding clusters in proportion to cumulative participant dropout cleared 80\% power in all six reported Table 3 scenarios. Dropout removes 15\% or 45\% of participants by the last occasion and 7.5\% or 22.5\% of observations across the four occasions under the design's schedule. Table 3 used 15, 22 and 29 clusters in the lower dropout arm and 19, 28 and 36 in the higher arm. The published power range was .83961--.90204 and the mean reconstructed range was .83938--.90192. At 19 clusters, ICC .05 and 45\% cumulative dropout, the source's .88187 fell below the central 95\% mask interval [.883690, .885691] from 10,000 masks, whose mean was .884912. In a separate saved diagnostic, dense and block generalized least squares (GLS) agreed at an adverse mask with power .881846. The source value is feasible in the lower tail, and the original mask is unknown. See \href{https://github.com/alectionary/vibe-math/blob/master/formalization/kupzyk-power-assumptions/results/full/oversampling-v2.csv}{\nolinkurl{results/full/oversampling-v2.csv}}, \href{https://github.com/alectionary/vibe-math/blob/master/formalization/kupzyk-power-assumptions/outputs/table3-gap-summary.csv}{\nolinkurl{outputs/table3-gap-summary.csv}}, \href{https://github.com/alectionary/vibe-math/blob/master/formalization/kupzyk-power-assumptions/outputs/table3-independent-check.csv}{\nolinkurl{outputs/table3-independent-check.csv}}, and \href{https://github.com/alectionary/vibe-math/blob/master/formalization/kupzyk-power-assumptions/outputs/table3-allocation-summary.csv}{\nolinkurl{outputs/table3-allocation-summary.csv}}.
\label{dis:table3-mask-tail}\disagreement{table3-mask-tail}

\subsection{Planning with partially observed participants}\label{sec:partial}
The power calculations use all available occasions, including measurements contributed before a participant drops out. \citet[pp.~94--95, 99--100]{kupzyk2011} also proposed in prose that planners add clusters in proportion to the fraction of \emph{observations} lost across the study. We applied that rule to the six Table 3 settings using observation losses of 7.5\% and 22.5\%. We multiplied the baseline cluster counts by one plus the observation-loss fraction and rounded to the nearest cluster or upward. Across the 12 setting-and-rounding rows, mean powers ranged from 81.18\% to 85.46\% over 500 masks per row. All means exceeded 80\% and were closer to it than those from the participant-loss rule in the same settings (\href{https://github.com/alectionary/vibe-math/blob/master/formalization/kupzyk-power-assumptions/results/full/observation-loss-oversampling.csv}{\nolinkurl{results/full/observation-loss-oversampling.csv}}).

A completers-only analysis using listwise deletion would discard earlier observations from participants who later drop out. Its power is outside the design evaluated here. In the separate completers illustration, Kupzyk adds 15\% to a target of 200 and obtains an enrollment of 230 \citep[p.~100]{kupzyk2011}. At 15\% loss of initial enrollees, 230 gives $230(.85)=195.5$ expected completers. The minimum integer enrollment to reach an expectation of 200 is $\lceil 200/.85\rceil=236$, with 200.6 expected completers (\href{https://github.com/alectionary/vibe-math/blob/master/formalization/kupzyk-power-assumptions/results/full/completers-arithmetic.json}{\nolinkurl{results/full/completers-arithmetic.json}}).
\label{dis:completer-inflation}\disagreement{completer-inflation}

\section{What we did ourselves and what we took from others}\label{sec:design}
The reconstruction uses the population model, factor settings, original comparison targets and printed size frequencies from \citet[Chapters III--V and Appendices A--D]{kupzyk2011}. We extracted the published results, reconstructed the generator and tests, performed independent covariance and fitting checks, and compared fresh estimates with the published figures and tables. The model has occasions within people within randomized clusters. It has fixed effects for time, treatment and their interaction, independent Gaussian random intercepts and slopes at the person and cluster levels, and Gaussian occasion errors. The interaction is .5 under the alternative and zero in the null checks.

The 54 initial settings cross ICC .05, .10 and .15, 0\%, 5\% and 15\% interval attrition, 50/50 or 70/30 allocation, and equal, normal-shaped or skewed cluster sizes. The ICC factor refers to the correlation of individual slopes within a cluster, and there are four occasions and 20 participants per cluster on average. The exemplary designs with unequal sizes use the source's printed integer frequency schedules (\href{https://github.com/alectionary/vibe-math/blob/master/formalization/kupzyk-power-assumptions/outputs/normsizes.csv}{\nolinkurl{outputs/normsizes.csv}} and \href{https://github.com/alectionary/vibe-math/blob/master/formalization/kupzyk-power-assumptions/outputs/skewsizes.csv}{\nolinkurl{outputs/skewsizes.csv}}). The empirical simulations draw new cluster sizes from a normal distribution with mean 20 and standard deviation 2, round them to integers and constrain them to 15 through 25. For skewed sizes, they apply two transformations in alternation to draws above 20 and constrain the final sizes to 15 through 22 (\href{https://github.com/alectionary/vibe-math/blob/master/formalization/kupzyk-power-assumptions/scripts/model.R#L6-L20}{\nolinkurl{scripts/model.R}}). Attrition removes a subset of the original participants at each interval without reentry. The empirical generator draws new outcomes and dropout masks. The exemplary calculation holds the variance components at their generating values. The initial cluster counts were 16, 22 and 28 as ICC increased, and the final exemplary grid used 13, 19 and 25. The six realized oversampling scenarios in Table 3 differ from the plans in Figure 5 (\href{https://github.com/alectionary/vibe-math/blob/master/formalization/kupzyk-power-assumptions/source-values/figure-5.csv}{\nolinkurl{source-values/figure-5.csv}}).

Empirical power estimates use restricted maximum likelihood (REML) with a two-sided .05 interaction test. The source used SAS/IML and PROC MIXED. The SAS/STAT 15.2 manual documents the PROC MIXED REML default \citep[PROC MIXED statement, METHOD option]{sasmixed}. The R reconstruction profiles the REML objective and checks the fitted interaction and its standard error against lme4 \citep{bates2015} and nlme \citep{nlme2026}, two R mixed-model implementations. A generator check compared 540 paired outputs and their random number states for the direct and reference implementations in R. The source's SAS calls \texttt{RANNOR(0)} and \texttt{RANUNI(0)} provide no fixed seeds for replay. See \href{https://github.com/alectionary/vibe-math/blob/master/formalization/kupzyk-power-assumptions/rerun/validate_generator.R}{\nolinkurl{rerun/validate_generator.R}}, \href{https://github.com/alectionary/vibe-math/blob/master/formalization/kupzyk-power-assumptions/rerun/validate_profile.R}{\nolinkurl{rerun/validate_profile.R}} and \href{https://github.com/alectionary/vibe-math/blob/master/formalization/kupzyk-power-assumptions/results/engine-comparison.csv}{\nolinkurl{results/engine-comparison.csv}}.

Exemplary power conditions on the observed design matrix \(X\) and generating covariance \(V\). With \(e\) selecting the treatment-by-time coefficient \(\delta\), the GLS variance is \(v=e^\mathsf{T}(X^\mathsf{T}V^{-1}X)^{-1}e\). The noncentrality of the one-numerator-degree-of-freedom \(F\) test is \(\delta^2/v\). We checked the variance by three calculations using dense, block and reduced representations of the covariance (\href{https://github.com/alectionary/vibe-math/blob/master/formalization/kupzyk-power-assumptions/results/full/fresh-gls-validation.csv}{\nolinkurl{results/full/fresh-gls-validation.csv}}). The dropout calculations average power across independently sampled masks. The source's exemplary estimate is conditional on its realized design.

The formulaic calculation isolates the effect of the denominator degrees of freedom. In balanced complete designs with \(J\) clusters, \(n\) people per cluster and \(t\) occasions, Optimal Design uses \(J-2\), while Kupzyk's between--within Equation 4.1 uses \(Jnt-J-2\). At the initial baseline this gives 1,262 degrees of freedom and power .877206 under the between--within calculation, consistent with the source's approximately .88. The even balanced Optimal Design crossings are 16, 22 and 28, and allowing odd counts yields 16, 21 and 27. Formulaic between--within crossings are 13, 19 and 25. At the final baseline, formulaic and exemplary powers can differ even at a common denominator rule because the formulaic and realized design calculations differ. See \href{https://github.com/alectionary/vibe-math/blob/master/formalization/kupzyk-power-assumptions/results/full/formulaic.csv}{\nolinkurl{results/full/formulaic.csv}} and \citet[pp.~83--89, 97, 103--104]{kupzyk2011}.

\section{How well did the power estimates reproduce?}\label{sec:grid}
The 54 initial empirical estimates all agreed at the scale of the adjusted comparison. We ran 25,600 new draws per cell, for 1,382,400 successful fits. The source lost one fit in each of four cells. Reading the plotted values adds uncertainty to both the comparison and the estimates of the published Monte Carlo errors. We combined the two binomial Monte Carlo standard errors in quadrature and added the figure-reading allowance to the resulting band. A Monte Carlo standard error for a rejection fraction \(\hat p\) from \(R\) independent draws is \(\sqrt{\hat p(1-\hat p)/R}\) \citep{morris2019}. Two initial empirical cells exceeded the ordinary band of two combined standard errors plus reading error, and both were inside the two-sided 95\% band with a Bonferroni adjustment across 54 cells \citep{dunn1961}. At ICC .10 with 70/30 allocation and normal sizes, the 5\% attrition estimates were .792557 versus .780898, and the 15\% estimates were .754458 versus .763672. The adjusted half-widths were about .0143 and .0141, respectively. These comparisons are checks at the resolution of the source and Monte Carlo estimates. See \citet[Figures 6--8]{kupzyk2011} and \href{https://github.com/alectionary/vibe-math/blob/master/formalization/kupzyk-power-assumptions/results/full/power-cell-comparison.csv}{\nolinkurl{results/full/power-cell-comparison.csv}}.

Exemplary power was also consistent at the resolution of the plotted values. We reconstructed all 54 initial and all 54 final settings. Three initial markers are hidden by other marks, leaving 51 readable initial comparisons. All 51 fit the adjusted bands, and all 54 final plotted comparisons fit their bands. The 18 no-dropout Table 1 values reproduce to five decimal places. With dropout, the variation among realized masks is distinct from the Monte Carlo error in estimating mean power over masks. A comparison of the source point with the mean uses the mask spread to account for the unknown source mask. See \href{https://github.com/alectionary/vibe-math/blob/master/formalization/kupzyk-power-assumptions/source-values/figures-6-11.csv}{\nolinkurl{source-values/figures-6-11.csv}}, \href{https://github.com/alectionary/vibe-math/blob/master/formalization/kupzyk-power-assumptions/source-values/table-1.csv}{\nolinkurl{source-values/table-1.csv}}, \href{https://github.com/alectionary/vibe-math/blob/master/formalization/kupzyk-power-assumptions/results/full/exemplary-v2.csv}{\nolinkurl{results/full/exemplary-v2.csv}} and \href{https://github.com/alectionary/vibe-math/blob/master/formalization/kupzyk-power-assumptions/results/full/power-cell-comparison.csv}{\nolinkurl{results/full/power-cell-comparison.csv}}.

Empirical power was usually below exemplary power. The source reports 41 of 54 initial cells below, and 39 of the 51 readable pairs are below. We find 40 of 54 below, with empirical minus exemplary power ranging from $-.01661$ to .00677. A separate 50,000-draw baseline run yielded .86842 against .877206 theoretical exemplary power, a deficit of .879 pp versus about .9 pp reported by \citet[pp.~106--107]{kupzyk2011}. See \href{https://github.com/alectionary/vibe-math/blob/master/formalization/kupzyk-power-assumptions/results/full/empirical-exemplary-gaps.csv}{\nolinkurl{results/full/empirical-exemplary-gaps.csv}} and \href{https://github.com/alectionary/vibe-math/blob/master/formalization/kupzyk-power-assumptions/results/full/empirical.csv}{\nolinkurl{results/full/empirical.csv}}.

\citet[pp.~106--107]{kupzyk2011} suggested that the gap would likely vanish with more simulation draws and that the two methods would then be equivalent. In a paired comparison on the same 50,000 balanced baseline outcomes, we held the reference critical value fixed and obtained power .87752 with the generating covariance and .86842 after refitting the covariance. The 43,421 refit rejections exactly reproduce the main follow-up count. Refitting lowered power by .00910 (paired Monte Carlo standard error .001068), supporting a covariance-refitting contribution to the baseline gap. Increasing the simulation count reduces Monte Carlo error around each rejection probability. This finite comparison is limited to the balanced baseline, and general convergence or equivalence remains unresolved (\href{https://github.com/alectionary/vibe-math/blob/master/formalization/kupzyk-power-assumptions/results/full/baseline-paired.csv}{\nolinkurl{results/full/baseline-paired.csv}}).
\label{dis:power-equivalence}\disagreement{power-equivalence}

The six Appendix C parameter tables agree at adjusted diagnostic precision. All 324 comparisons lie within their Bonferroni bands when the source's unavailable Monte Carlo standard errors are approximated by ours. The six parameters are the person and cluster intercept and slope variances, the slope ICC, and the time-by-treatment interaction. Boundary estimates, where a fitted variance equals zero, are distinct from failed fits. The initial baseline boundary fraction in R was .07266. The slope ICC is computed within each fit before averaging, so the average of the fitted ratios can differ from a ratio of averaged variances. See \href{https://github.com/alectionary/vibe-math/blob/master/formalization/kupzyk-power-assumptions/results/full/appendix-c-comparison.csv}{\nolinkurl{results/full/appendix-c-comparison.csv}} and \href{https://github.com/alectionary/vibe-math/blob/master/formalization/kupzyk-power-assumptions/results/full/empirical.csv}{\nolinkurl{results/full/empirical.csv}}.

\section{Null calibration and denominator rules}\label{sec:null}
The two Appendix D settings show that denominator rules can change both power and false positive rates. We used 25,600 draws under the null per setting. The corresponding alternative rates use the initial grid draws. We compared Optimal Design's $J-2$ rule, between--within, containment, Satterthwaite and Kenward--Roger. The SAS manual documents the between--within and containment rules \citep[MODEL statement, DDFM options]{sasmixed}. We checked our Satterthwaite calculation \citep{satterthwaite1946} against the lmerTest implementation \citep{kuznetsova2017}. Kenward--Roger also adjusts the covariance of the fixed effects. Its implementation was checked against pbkrtest \citep{halekoh2014} (\href{https://github.com/alectionary/vibe-math/blob/master/formalization/kupzyk-power-assumptions/results/full/kr-validation.csv}{\nolinkurl{results/full/kr-validation.csv}}).

At baseline, between--within has a Type I error rate of .06500, compared with .06839 in Table D.1. Both exceed the nominal .05, as does containment. The reconstructed Kenward--Roger rate is .04625 with all variance components retained and .04891 when zero components are dropped. The source reports .05021. The reconstructed and source rates for all five rules fall within Bradley’s liberal robustness interval [.025, .075] around .05 \citep{bradley1978}, as applied by \citet[p.~105]{kupzyk2011}. The source's own two-standard-error interval for Satterthwaite and Kenward--Roger, [.04744, .05279], centers on .050115 rather than the reported .05021. The .000095 offset exceeds the .00001 rounding allowance. If the interval was symmetric at two standard errors, using the printed point in the variance formula implies about 26,658 draws, versus about 25,000 from most neighboring widths. These calculations leave the original interval procedure unresolved (\href{https://github.com/alectionary/vibe-math/blob/master/formalization/kupzyk-power-assumptions/results/appendix-d-source-identity.csv}{\nolinkurl{results/appendix-d-source-identity.csv}}).
\label{dis:appendix-d1-interval}\disagreement{appendix-d1-interval}

The point classifications differ under the narrower robustness interval [.035, .065] that \citet[p.~105]{kupzyk2011} attributes to \citet{serlin2000}. The published between--within and containment rates of .06839 and .06811 exceed .065. The corresponding R estimates, .06500 and .0648828, lie at and just inside that endpoint. Both R rates are within .00012 of the cutoff, with Monte Carlo standard errors of about .00154, so the change in classification at the point estimates does not establish a difference in the underlying classification (\href{https://github.com/alectionary/vibe-math/blob/master/formalization/kupzyk-power-assumptions/results/full/appendix-d-comparison.csv}{\nolinkurl{results/full/appendix-d-comparison.csv}}).
\label{dis:appendix-d1-serlin}\disagreement{appendix-d1-serlin}

At ICC .15 with skewed sizes, 15\% interval attrition, and 70/30 allocation, the $J-2$ rule differs materially from Table D.2. Its power and Type I error at $J=28$ and 26 degrees of freedom were .69594 and .04969, versus .67128 and .04279 in the source. The rerun's two-standard-error bands are approximately [.69019, .70169] and [.04697, .05240]. The published null interval [.04023, .04535] excludes the rerun point. A replay recovered 17,816 alternative and 1,272 null rejections at 26 degrees of freedom. Recomputing the tests on the same fits at 15 degrees of freedom gives .66953 and .04270. This sensitivity explains how a smaller effective denominator could produce both rates. The original denominator and fitted state are unresolved. See \href{https://github.com/alectionary/vibe-math/blob/master/formalization/kupzyk-power-assumptions/results/full/appendix-d-comparison.csv}{\nolinkurl{results/full/appendix-d-comparison.csv}}, \href{https://github.com/alectionary/vibe-math/blob/master/formalization/kupzyk-power-assumptions/results/full/appendix-d2-df-sweep.csv}{\nolinkurl{results/full/appendix-d2-df-sweep.csv}} and \href{https://github.com/alectionary/vibe-math/blob/master/formalization/kupzyk-power-assumptions/results/full/kenward-roger-summary.csv}{\nolinkurl{results/full/kenward-roger-summary.csv}}.
\label{dis:appendix-d2-od}\disagreement{appendix-d2-od}

\section{Worked examples and simulation precision}\label{sec:examples}
The G*Power example reproduces to six decimals, and the two Optimal Design display values differ from the reconstructed formula. The published and recalculated G*Power values are both .584382. For the two Optimal Design examples they are .550898 versus .553630 and .422156 versus .425064. Computing the tail using either noncentral $t$ or its square as a noncentral $F$ gives the same recalculated values. The two differences, about .003 each, leave open the effect of a software convention or a different effective input. See \citet[Figures 1--3]{kupzyk2011}, \href{https://github.com/alectionary/vibe-math/blob/master/formalization/kupzyk-power-assumptions/source-values/figures-2-3.csv}{\nolinkurl{source-values/figures-2-3.csv}} and \href{https://github.com/alectionary/vibe-math/blob/master/formalization/kupzyk-power-assumptions/results/source-identities.json}{\nolinkurl{results/source-identities.json}}.
\label{dis:od-screen-example}\disagreement{od-screen-example}

The simulation-precision calculation supports the recommended draw counts. At power .80, a normal approximation with two standard errors on either side gives full width \(4\sqrt{.8(.2)/R}\), so widths .01 and .005 require 25,600 and 102,400 draws. At 1,000 draws, the exact central binomial quantiles expressed as proportions are .775 and .824. Figure 4 shows two 500-draw illustrations. Our independent illustration uses a fixed standardized predictor and calibrates its population power to .80 at $n=75$. Across 10,000 simulated power curves at each draw count, first 80\% crossings ranged from 71 to 80 at 500 draws and 74 to 77 at 25,600 draws. These ranges illustrate Monte Carlo variability. The source toy model and its two random streams are unavailable. See \citet[pp.~59--64]{kupzyk2011}, \href{https://github.com/alectionary/vibe-math/blob/master/formalization/kupzyk-power-assumptions/source-values/figure-4.csv}{\nolinkurl{source-values/figure-4.csv}} and \href{https://github.com/alectionary/vibe-math/blob/master/formalization/kupzyk-power-assumptions/results/figure4-independent-summary.json}{\nolinkurl{results/figure4-independent-summary.json}}.

\section{Assumptions}\label{sec:assumptions}

\subsection{Observed-count degrees of freedom}\label{sec:observed}
The between--within implementation extends the complete-design rule to $M-J-2$, where $M$ is the observed number of occasions summed over participants. This treats lost observations as reducing the within count. It equals Equation 4.1 for a complete balanced design. The extension is a reconstruction choice for the missing data settings (\href{https://github.com/alectionary/vibe-math/blob/master/formalization/kupzyk-power-assumptions/results/design-checks.csv}{\nolinkurl{results/design-checks.csv}}).

\subsection{Independent dropout realizations}\label{sec:masks}
The mask average estimates power under the reported sampling rule with independent realizations of missingness. The source's exemplary powers, oversampling cells, and cluster crossings are conditional on realized designs. The two estimands can differ even when they use identical covariance and test formulas. We retain the mask distribution to distinguish near-target threshold differences and lower-tail comparisons from fitting errors.

\subsection{Unobserved source Monte Carlo variance}
Published Appendix C summaries give means without the draw-level spread needed to compute their Monte Carlo errors. We use the corresponding rerun standard error as a proxy. The adjusted diagnostic bands depend on that approximation.

\section{Ambiguities that were resolved}\label{sec:ambiguities}

\subsection{The Table 4 high-attrition column}
Table 4 prints 10\% per interval in the second attrition column. The inspected page confirms the label, and Figure 5 specifies 0\%, 5\% and 15\%. Recomputing the 0-minus-15\% contrasts from Table 1 gives 6.1692, 3.8722 and 3.1325 pp, which round to the column's 6.2, 3.9 and 3.1. We treat this as a source label error. It does not change the numerical findings (\href{https://github.com/alectionary/vibe-math/blob/master/formalization/kupzyk-power-assumptions/source-values/table-4.csv}{\nolinkurl{source-values/table-4.csv}} and \href{https://github.com/alectionary/vibe-math/blob/master/formalization/kupzyk-power-assumptions/results/source-identities.json}{\nolinkurl{results/source-identities.json}}).
\label{dis:table4-high-attrition-label}\disagreement{table4-high-attrition-label}

\subsection{The oversampling rule and the two reference counts}
The printed oversampling counts follow the final baseline counts multiplied by one plus cumulative dropout and rounded to the nearest cluster. The proposed counts in Figure 5 follow a different grid starting at the initial Optimal Design counts. We compared power at the six realized counts in Table 3 and evaluated the four research questions at the final exemplary counts in Table 1 (\href{https://github.com/alectionary/vibe-math/blob/master/formalization/kupzyk-power-assumptions/results/source-identities.json}{\nolinkurl{results/source-identities.json}}).

\subsection{Hidden markers and final figure values}
Three exemplary markers in the initial figures are occluded. We left these observations missing in figure comparisons and used the independently reported full-grid count for the empirical-exemplary sign comparison. Table 1 provides the final numerical comparison and calibrates the reading of Figures 9--11 (\href{https://github.com/alectionary/vibe-math/blob/master/formalization/kupzyk-power-assumptions/results/figure-table1-check.csv}{\nolinkurl{results/figure-table1-check.csv}}).

\section{What remains unresolved}\label{sec:unresolved}
The historical inputs needed for exact replay are the fitted SAS records, random states, and realized designs. The SAS syntax and generator checks fix the modeled distributions. Agreement across R fitting and GLS implementations constrains errors in our reconstruction. The published summaries leave the following comparisons unresolved at the scale indicated.

\begin{itemize}
\item For the three threshold differences, the neighboring-count scan and independent covariance check support the rerun crossings. Mask variation is the leading explanation. The realized masks at the published counts and their predecessors would allow a direct crossing comparison.
\item For the 19-cluster Table 3 cell, an unusually adverse mask is feasible under the reconstructed rule. The allocation diagnostic and covariance cross-check do not identify the source configuration. The realized design and mask would distinguish sampling variation from allocation or rounding conventions.
\item For Table D.2, a smaller effective denominator brings both rates near the printed pair and leads the explanations tested. Replay rules out a counting discrepancy between the R paths. The original per-fit test statistics, denominator degrees of freedom and rejection counts would distinguish a denominator convention from a different fitted distribution.
\item For the Table D.1 intervals, transcription and ordinary rounding have been checked. The remaining candidates include a different rate, draw count, or interval procedure, with insufficient evidence to order them. A table-generation record with the original counts and interval formula would settle the calculation.
\item For the two Optimal Design screens, equivalent tail formulas reproduce the same gaps. A software convention or a different effective input remains plausible. The original executable or saved project with effective inputs would distinguish these explanations. The Figure 4 seeds and generator would permit comparison with the exact plotted toy realizations.
\end{itemize}

\section{Code and source availability}\label{sec:availability}
The University of Nebraska--Lincoln repository record identifies the dissertation and its canonical download. A complete archived copy of the university export was obtained on 2026-11-09. Its SHA-256 is \nolinkurl{368dc575d9f744bf5ca5516873a633691d999fbdd5d19516994c64e10674896a}. It has 143 physical pages, with a ten-page offset from the printed Arabic numbering. The institutional record, direct PDF response and successful archival route are documented in \href{https://github.com/alectionary/vibe-math/blob/master/formalization/kupzyk-power-assumptions/provenance/source-acquisition.json}{\nolinkurl{provenance/source-acquisition.json}}. The copyright-restricted PDF is held separately. The source transcriptions and figure measurements are in the repository. See \href{https://github.com/alectionary/vibe-math/blob/master/formalization/kupzyk-power-assumptions/source-inventory.json}{\nolinkurl{source-inventory.json}}.

SAS/IML, SAS PROC MIXED and Optimal Design 2.0 were unavailable for execution in this study. The engine record describes the R reconstruction and the loss of access to SAS-specific behavior and Optimal Design output (\href{https://github.com/alectionary/vibe-math/blob/master/formalization/kupzyk-power-assumptions/notes/rerun-engine.json}{\nolinkurl{notes/rerun-engine.json}}). The R and Python code and saved result tables are linked in Appendix~\ref{sec:artifacts}.

\section{Next steps}\label{sec:next}
Original design and fit records would turn several sensitivity comparisons into replay checks. A unit-level design and mask would let the same GLS calculation reproduce the published exemplary estimand at Table 3 and the three Table 5 thresholds. Per-fit test statistics and denominator degrees of freedom from Appendix D would let us recompute the extreme rates under the actual rule. An archived Optimal Design project would allow the reported and reconstructed effective inputs to be compared directly.

\section{Numerical comparison at printed precision}\label{sec:numbers}
The table distinguishes numerical agreement with a printed value from reconstruction scope. At the indicated precision, \emph{match} means our value rounds to the published value, \emph{close} means the gap is at most one unit in the last printed place, and \emph{mismatch} means the gap exceeds that unit. All rows have limited scope because they use independently calculated powers or newly simulated data. The numerical labels use rounding, with Monte Carlo and mask comparisons reported separately above.

\begin{longtable}{@{}L{.27\linewidth}rrL{.14\linewidth}ll@{}}
\caption{Selected printed estimates. Sources are Kupzyk's Figures 2--3, Tables 1, 3 and 5 and Appendix D, and the corresponding reconstructed values.}\\
\toprule
Claim & Published & Ours & Precision & Numerical & Scope \\
\midrule\endfirsthead
\toprule
Claim & Published & Ours & Precision & Numerical & Scope \\
\midrule\endhead
G*Power example & .584382 & .584382 & 6 decimals & match & limited \\
Optimal Design example, ICC 0 & .550898 & .553630 & 6 decimals & mismatch & limited \\
Optimal Design example, ICC .10 & .422156 & .425064 & 6 decimals & mismatch & limited \\
Final baseline exemplary power & .80081 & .80081 & 5 decimals & match & limited \\
Table 3, 19 clusters & .88187 & .88491 & 5 decimals & mismatch & limited \\
Threshold, cell 11 & 17 & 16 & 1 cluster & close & limited \\
Threshold, cell 13 & 16 & 15 & 1 cluster & close & limited \\
Threshold, cell 54 & 33 & 32 & 1 cluster & close & limited \\
D.1 between--within Type I error & .06839 & .06500 & 5 decimals & mismatch & limited \\
D.1 Kenward--Roger Type I error & .05021 & .04625 & 5 decimals & mismatch & limited \\
D.2 Optimal Design power & .67128 & .69594 & 5 decimals & mismatch & limited \\
D.2 Optimal Design Type I error & .04279 & .04969 & 5 decimals & mismatch & limited \\
\bottomrule
\end{longtable}

The worked example matches for the noncentral $t$ calculation and differs for the two software displays. At the final baseline the fixed design yields the same printed power. At 19 clusters the source value is in the far lower tail of the reconstructed mask distribution. All three listed thresholds differ by exactly one cluster. The two D.1 rate gaps exceed the rounding unit. All baseline rates lie inside the broad Bradley practical band [.025, .075]. The between--within and containment point estimates change classification under the narrower Serlin band [.035, .065], with the R values on or just inside .065. Both D.2 gaps exceed the rounding unit and the rerun two-standard-error bands.

\appendix
\section{Artifacts and regeneration}\label{sec:artifacts}
Use the repository's \texttt{master} branch. Set the study path from the repository root before using either regeneration route.
\begin{quote}
\small\verb|study=formalization/kupzyk-power-assumptions|
\end{quote}

\begin{description}
\item[R, numerical reconstruction.] Run
\begin{quote}
\small\verb|sh "$study/rerun/reproduce.sh" /tmp/kupzyk-reproduction 4|
\end{quote}
\href{https://www.r-project.org/}{R 4.5.3} supplied random draws, distribution functions and numerical optimization \citep{rcore2026}. \href{https://cran.r-project.org/package=Rcpp}{Rcpp 1.1.2} compiled the C++ profile evaluator \citep{eddelbuettel2011}. The fitting and denominator checks used \href{https://cran.r-project.org/package=lme4}{lme4 2.0-6} \citep{bates2015}, \href{https://cran.r-project.org/package=nlme}{nlme 3.1-170} \citep{nlme2026}, \href{https://cran.r-project.org/package=lmerTest}{lmerTest 3.2-1} \citep{kuznetsova2017}, \href{https://cran.r-project.org/package=pbkrtest}{pbkrtest 0.5.5} \citep{halekoh2014}, and \href{https://cran.r-project.org/package=Matrix}{Matrix 1.7-5} for rank calculations \citep{matrix2026}.
\href{https://www.python.org/}{Python 3.14.10} aggregated and compared results \citep{python2026}. The independent precision calculation used \href{https://numpy.org/}{NumPy 2.5.4} for arrays \citep{harris2020} and \href{https://scipy.org/}{SciPy 1.18.1} for numerical roots and distributions \citep{virtanen2020}. \href{https://pyyaml.org/}{PyYAML 6.0.3} parsed the evidence records \citep{pyyaml2026}. \href{https://github.com/pymupdf/PyMuPDF}{PyMuPDF 2.1} extracted source-figure geometry, and \href{https://github.com/pymupdf/PyMuPDF}{PyMuPDF 1.26.6} extracted annotations from the replication PDF for link review \citep{pymupdf2026}.
The environment is recorded in \href{https://github.com/alectionary/vibe-math/blob/master/formalization/kupzyk-power-assumptions/provenance/rerun-runtime.txt}{\nolinkurl{provenance/rerun-runtime.txt}}. See the prerequisites and checkpoints in \href{https://github.com/alectionary/vibe-math/blob/master/formalization/kupzyk-power-assumptions/rerun/README.md}{\nolinkurl{rerun/README.md}} and the exact command chain in \href{https://github.com/alectionary/vibe-math/blob/master/formalization/kupzyk-power-assumptions/rerun/reproduce.sh}{\nolinkurl{rerun/reproduce.sh}}. The full empirical grid schedules 128 batches of 200 draws per setting, plus the 50,000-draw follow-up and two null cases.  The following paths are relative to the selected run directory. Each link opens the corresponding checked-in result.

\begin{itemize}
\item \href{https://github.com/alectionary/vibe-math/blob/master/formalization/kupzyk-power-assumptions/rerun/formulaic.R}{\nolinkurl{rerun/formulaic.R}} writes \href{https://github.com/alectionary/vibe-math/blob/master/formalization/kupzyk-power-assumptions/results/full/formulaic.csv}{\nolinkurl{formulaic.csv}}.
\item \href{https://github.com/alectionary/vibe-math/blob/master/formalization/kupzyk-power-assumptions/rerun/precision_demo.py}{\nolinkurl{rerun/precision_demo.py}} writes \href{https://github.com/alectionary/vibe-math/blob/master/formalization/kupzyk-power-assumptions/results/figure4-independent-summary.json}{\nolinkurl{figure4-independent-summary.json}}.
\item \href{https://github.com/alectionary/vibe-math/blob/master/formalization/kupzyk-power-assumptions/rerun/validate_profile.R}{\nolinkurl{rerun/validate_profile.R}} writes \href{https://github.com/alectionary/vibe-math/blob/master/formalization/kupzyk-power-assumptions/results/engine-comparison.csv}{\nolinkurl{engine-comparison.csv}}.
\item \href{https://github.com/alectionary/vibe-math/blob/master/formalization/kupzyk-power-assumptions/rerun/check_oversampling.py}{\nolinkurl{rerun/check_oversampling.py}} writes \href{https://github.com/alectionary/vibe-math/blob/master/formalization/kupzyk-power-assumptions/results/full/observation-loss-oversampling.csv}{\nolinkurl{comparisons/observation-loss-oversampling.csv}} and \href{https://github.com/alectionary/vibe-math/blob/master/formalization/kupzyk-power-assumptions/results/full/completers-arithmetic.json}{\nolinkurl{comparisons/completers-arithmetic.json}}.
\item \href{https://github.com/alectionary/vibe-math/blob/master/formalization/kupzyk-power-assumptions/rerun/baseline_paired.R}{\nolinkurl{rerun/baseline_paired.R}} writes \href{https://github.com/alectionary/vibe-math/blob/master/formalization/kupzyk-power-assumptions/results/full/baseline-paired.csv}{\nolinkurl{baseline-paired.csv}}.
\end{itemize}

Reuse the checkpoint directory for a resumed identical run. Use a new directory after changes to scientific inputs, versions or seeds.
\item[S, source measurements.] Retrieve the archived PDF recorded in \href{https://github.com/alectionary/vibe-math/blob/master/formalization/kupzyk-power-assumptions/provenance/source-acquisition.json}{\nolinkurl{provenance/source-acquisition.json}} and verify its SHA-256. \href{https://poppler.freedesktop.org/}{Poppler 26.10.0} extracts layout text from the source \citep{poppler2026}. Use the source audit instructions in \href{https://github.com/alectionary/vibe-math/blob/master/formalization/kupzyk-power-assumptions/rerun/README.md}{\nolinkurl{rerun/README.md}}. This step regenerates the source transcriptions and calculations from the archival PDF. Route R uses the checked-in transcriptions and can run without the archival file.
\end{description}

The report rebuild uses the saved results. With \href{https://github.com/tectonic-typesetting/tectonic}{Tectonic 0.17.0} \citep{tectonic2026} and its LaTeX bundle, run the following command from the repository root.
\begin{quote}
\small\verb|(cd "$study/paper" && tectonic kupzyk-replication.tex)|
\end{quote}
The bibliography is embedded in the editable source. This rebuild writes the PDF alongside the source.

\small
\begin{longtable}{@{}L{.28\linewidth}L{.60\linewidth}l@{}}
\caption{Evidence by claim and regeneration route. Published page locators refer to the printed dissertation.}\\
\toprule
Claim & Evidence & Route \\
\midrule\endfirsthead
\toprule
Claim & Evidence & Route \\
\midrule\endhead
Chapter I motivation (pp. 1--6) & \href{https://github.com/alectionary/vibe-math/blob/master/formalization/kupzyk-power-assumptions/notes/claim-inventory.md}{\nolinkurl{notes/claim-inventory.md}} & S \\
Chapter II software examples & \href{https://github.com/alectionary/vibe-math/blob/master/formalization/kupzyk-power-assumptions/results/source-identities.json}{\nolinkurl{results/source-identities.json}} & S \\
Figure 4 simulation precision & \href{https://github.com/alectionary/vibe-math/blob/master/formalization/kupzyk-power-assumptions/results/figure4-independent-summary.json}{\nolinkurl{results/figure4-independent-summary.json}} & R \\
Chapter III crossed design & \href{https://github.com/alectionary/vibe-math/blob/master/formalization/kupzyk-power-assumptions/source-values/figure-5.csv}{\nolinkurl{source-values/figure-5.csv}} & S \\
Initial empirical grid & \href{https://github.com/alectionary/vibe-math/blob/master/formalization/kupzyk-power-assumptions/results/full/empirical.csv}{\nolinkurl{results/full/empirical.csv}} & R \\
Initial exemplary grid & \href{https://github.com/alectionary/vibe-math/blob/master/formalization/kupzyk-power-assumptions/results/full/exemplary-v2.csv}{\nolinkurl{results/full/exemplary-v2.csv}} & R \\
Final exemplary grid & \href{https://github.com/alectionary/vibe-math/blob/master/formalization/kupzyk-power-assumptions/results/full/power-cell-comparison.csv}{\nolinkurl{results/full/power-cell-comparison.csv}} & R \\
Optimal Design formula & \href{https://github.com/alectionary/vibe-math/blob/master/formalization/kupzyk-power-assumptions/results/full/formulaic.csv}{\nolinkurl{results/full/formulaic.csv}} & R \\
Between--within formula & \href{https://github.com/alectionary/vibe-math/blob/master/formalization/kupzyk-power-assumptions/results/full/fresh-gls-validation.csv}{\nolinkurl{results/full/fresh-gls-validation.csv}} & R \\
Chapter IV results & \href{https://github.com/alectionary/vibe-math/blob/master/formalization/kupzyk-power-assumptions/results/full/comparison-summary.json}{\nolinkurl{results/full/comparison-summary.json}} & R \\
Attrition contrasts & \href{https://github.com/alectionary/vibe-math/blob/master/formalization/kupzyk-power-assumptions/results/full/paired-contrasts.csv}{\nolinkurl{results/full/paired-contrasts.csv}} & R \\
Table 3 and observation-loss oversampling & \href{https://github.com/alectionary/vibe-math/blob/master/formalization/kupzyk-power-assumptions/results/full/oversampling-v2.csv}{\nolinkurl{results/full/oversampling-v2.csv}}, \href{https://github.com/alectionary/vibe-math/blob/master/formalization/kupzyk-power-assumptions/results/full/observation-loss-oversampling.csv}{\nolinkurl{results/full/observation-loss-oversampling.csv}} & R \\
Allocation contrast & \href{https://github.com/alectionary/vibe-math/blob/master/formalization/kupzyk-power-assumptions/results/full/paired-contrasts.csv}{\nolinkurl{results/full/paired-contrasts.csv}} & R \\
Size contrasts & \href{https://github.com/alectionary/vibe-math/blob/master/formalization/kupzyk-power-assumptions/results/full/paired-contrasts.csv}{\nolinkurl{results/full/paired-contrasts.csv}} & R \\
ICC moderation & \href{https://github.com/alectionary/vibe-math/blob/master/formalization/kupzyk-power-assumptions/results/full/comparison-summary.json}{\nolinkurl{results/full/comparison-summary.json}} & R \\
Predicted loss ordering & \href{https://github.com/alectionary/vibe-math/blob/master/formalization/kupzyk-power-assumptions/source-values/table-1.csv}{\nolinkurl{source-values/table-1.csv}}, \href{https://github.com/alectionary/vibe-math/blob/master/formalization/kupzyk-power-assumptions/results/full/paired-contrasts.csv}{\nolinkurl{results/full/paired-contrasts.csv}} & S, R \\
Table 5 thresholds & \href{https://github.com/alectionary/vibe-math/blob/master/formalization/kupzyk-power-assumptions/results/full/threshold-comparison.csv}{\nolinkurl{results/full/threshold-comparison.csv}} & R \\
Chapter V interpretation & \href{https://github.com/alectionary/vibe-math/blob/master/formalization/kupzyk-power-assumptions/results/full/appendix-d-comparison.csv}{\nolinkurl{results/full/appendix-d-comparison.csv}}, \href{https://github.com/alectionary/vibe-math/blob/master/formalization/kupzyk-power-assumptions/results/full/completers-arithmetic.json}{\nolinkurl{results/full/completers-arithmetic.json}} & R \\
50,000-draw baseline & \href{https://github.com/alectionary/vibe-math/blob/master/formalization/kupzyk-power-assumptions/results/full/empirical.csv}{\nolinkurl{results/full/empirical.csv}}, \href{https://github.com/alectionary/vibe-math/blob/master/formalization/kupzyk-power-assumptions/results/full/baseline-paired.csv}{\nolinkurl{results/full/baseline-paired.csv}} & R \\
Empirical--exemplary gap & \href{https://github.com/alectionary/vibe-math/blob/master/formalization/kupzyk-power-assumptions/results/full/empirical-exemplary-gaps.csv}{\nolinkurl{results/full/empirical-exemplary-gaps.csv}} & R \\
Appendix A generator and REML & \href{https://github.com/alectionary/vibe-math/blob/master/formalization/kupzyk-power-assumptions/rerun/validate_generator.R}{\nolinkurl{rerun/validate_generator.R}}, \href{https://github.com/alectionary/vibe-math/blob/master/formalization/kupzyk-power-assumptions/results/engine-comparison.csv}{\nolinkurl{results/engine-comparison.csv}} & R \\
Appendix B covariance and sizes & \href{https://github.com/alectionary/vibe-math/blob/master/formalization/kupzyk-power-assumptions/results/full/fresh-gls-validation.csv}{\nolinkurl{results/full/fresh-gls-validation.csv}}, \href{https://github.com/alectionary/vibe-math/blob/master/formalization/kupzyk-power-assumptions/outputs/normsizes.csv}{\nolinkurl{outputs/normsizes.csv}}, \href{https://github.com/alectionary/vibe-math/blob/master/formalization/kupzyk-power-assumptions/outputs/skewsizes.csv}{\nolinkurl{outputs/skewsizes.csv}} & R, S \\
Appendix C fitted means & \href{https://github.com/alectionary/vibe-math/blob/master/formalization/kupzyk-power-assumptions/results/full/appendix-c-comparison.csv}{\nolinkurl{results/full/appendix-c-comparison.csv}} & R \\
Appendix D.1 baseline & \href{https://github.com/alectionary/vibe-math/blob/master/formalization/kupzyk-power-assumptions/results/full/kenward-roger-summary.csv}{\nolinkurl{results/full/kenward-roger-summary.csv}} & R \\
Appendix D.2 extreme setting & \href{https://github.com/alectionary/vibe-math/blob/master/formalization/kupzyk-power-assumptions/results/full/appendix-d2-df-sweep.csv}{\nolinkurl{results/full/appendix-d2-df-sweep.csv}} & R \\
Practical planning ordering & \href{https://github.com/alectionary/vibe-math/blob/master/formalization/kupzyk-power-assumptions/results/full/paired-contrasts.csv}{\nolinkurl{results/full/paired-contrasts.csv}}, \href{https://github.com/alectionary/vibe-math/blob/master/formalization/kupzyk-power-assumptions/results/full/threshold-comparison.csv}{\nolinkurl{results/full/threshold-comparison.csv}} & R \\
\bottomrule
\end{longtable}
\normalsize

\section{Compliance ledger}\label{sec:ledger}
\begin{tabular}{@{}L{.24\linewidth}L{.16\linewidth}L{.55\linewidth}@{}}
\toprule
Requirement & Status & Evidence \\
\midrule
Source verification & followed & \href{https://github.com/alectionary/vibe-math/blob/master/formalization/kupzyk-power-assumptions/provenance/source-acquisition.json}{\nolinkurl{provenance/source-acquisition.json}} records the complete archival copy. \\
Design coverage & followed & \href{https://github.com/alectionary/vibe-math/blob/master/formalization/kupzyk-power-assumptions/REPLICATION-PLAN.yaml}{\nolinkurl{REPLICATION-PLAN.yaml}} and \href{https://github.com/alectionary/vibe-math/blob/master/formalization/kupzyk-power-assumptions/results/full/comparison-summary.json}{\nolinkurl{results/full/comparison-summary.json}} account for the crossed grid and follow-ups. \\
Original engine & departed & SAS and Optimal Design were unavailable. The independent R checks and limits are described in \href{https://github.com/alectionary/vibe-math/blob/master/formalization/kupzyk-power-assumptions/notes/rerun-engine.json}{\nolinkurl{notes/rerun-engine.json}}. \\
Exact source streams & departed & The published code supplies zero-valued seed arguments and the unit-level masks are unavailable. New streams cover all settings. \\
Comparisons and uncertainty & followed & \href{https://github.com/alectionary/vibe-math/blob/master/formalization/kupzyk-power-assumptions/results/full/power-cell-comparison.csv}{\nolinkurl{results/full/power-cell-comparison.csv}} separates reading allowances, Monte Carlo error and mask spread. \\
Citation scope & followed & The \href{https://github.com/alectionary/vibe-math/blob/master/formalization/kupzyk-power-assumptions/citation-inventory.json}{\nolinkurl{citation-inventory.json}} maps methods and software to their citations and accounts for all earlier references. \\
Discrepancy record & followed & \href{https://github.com/alectionary/vibe-math/blob/master/formalization/kupzyk-power-assumptions/DISAGREEMENTS.yaml}{\nolinkurl{DISAGREEMENTS.yaml}} records the checks and unresolved cases. \\
\bottomrule
\end{tabular}

\begin{thebibliography}{99}
\bibitem[Bates et~al.(2015)Bates, Mächler, Bolker, and Walker]{bates2015}
Bates, D., Mächler, M., Bolker, B., and Walker, S. (2015).
Fitting linear mixed-effects models using lme4.
\emph{Journal of Statistical Software}, 67(1).
\href{https://www.jstatsoft.org/article/view/v067i01}{doi:10.18637/jss.v067.i01}.

\bibitem[Bates et~al.(2026)]{matrix2026}
Bates, D., Maechler, M., and Jagan, M. (2026).
\emph{Matrix: Sparse and dense matrix classes and methods}, [R package 1.7-5].
\href{https://doi.org/10.32614/CRAN.package.Matrix}{doi:10.32614/CRAN.package.Matrix}.

\bibitem[Bradley(1978)]{bradley1978}
Bradley, J. V. (1978). Robustness?
\emph{British Journal of Mathematical and Statistical Psychology}, 31(2), 144--152.
\href{https://doi.org/10.1111/j.2044-8317.1978.tb00581.x}{doi:10.1111/j.2044-8317.1978.tb00581.x}.

\bibitem[Dunn(1961)]{dunn1961}
Dunn, O. J. (1961). Multiple comparisons among means.
\emph{Journal of the American Statistical Association}, 56(293), 52--64.
\href{https://doi.org/10.1080/01621459.1961.10482090}{doi:10.1080/01621459.1961.10482090}.

\bibitem[Eddelbuettel and François(2011)]{eddelbuettel2011}
Eddelbuettel, D. and François, R. (2011). Rcpp: Seamless R and C++ integration.
\emph{Journal of Statistical Software}, 40(8), 1--18.
\href{https://doi.org/10.18637/jss.v040.i08}{doi:10.18637/jss.v040.i08}.

\bibitem[Halekoh and Højsgaard(2014)]{halekoh2014}
Halekoh, U. and Højsgaard, S. (2014).
A Kenward--Roger approximation and parametric bootstrap methods for tests in linear mixed models: The R package pbkrtest.
\emph{Journal of Statistical Software}, 59(9).
\href{https://www.jstatsoft.org/article/view/v059i09}{doi:10.18637/jss.v059.i09}.

\bibitem[Harris et~al.(2020)]{harris2020}
Harris, C. R., Millman, K. J., van der Walt, S. J., et al. (2020). Array programming with NumPy.
\emph{Nature}, 585, 357--362.
\href{https://doi.org/10.1038/s41586-020-2649-2}{doi:10.1038/s41586-020-2649-2}.

\bibitem[Kupzyk(2011)]{kupzyk2011}
Kupzyk, K. A. (2011).
\emph{The Effects of Simplifying Assumptions in Power Analysis}.
PhD dissertation, University of Nebraska--Lincoln.
\href{https://digitalcommons.unl.edu/cehsdiss/106/}{DigitalCommons@University of Nebraska--Lincoln, record 106}.

\bibitem[Kuznetsova et~al.(2017)]{kuznetsova2017}
Kuznetsova, A., Brockhoff, P. B., and Christensen, R. H. B. (2017). lmerTest package: Tests in linear mixed effects models.
\emph{Journal of Statistical Software}, 82(13), 1--26.
\href{https://doi.org/10.18637/jss.v082.i13}{doi:10.18637/jss.v082.i13}.

\bibitem[Morris et~al.(2019)Morris, White, and Crowther]{morris2019}
Morris, T. P., White, I. R., and Crowther, M. J. (2019).
Using simulation studies to evaluate statistical methods.
\emph{Statistics in Medicine}, 38(11), 2074--2102.
\href{https://doi.org/10.1002/sim.8086}{doi:10.1002/sim.8086}.
\bibitem[Pinheiro et~al.(2026)]{nlme2026}
Pinheiro, J., Bates, D., and R Core Team. (2026).
\emph{nlme: Linear and nonlinear mixed effects models}, [R package 3.1-170].
\href{https://doi.org/10.32614/CRAN.package.nlme}{doi:10.32614/CRAN.package.nlme}.

\bibitem[Poppler Project(2026)]{poppler2026}
Poppler Project. (2026).
\emph{Poppler}, [PDF library and tools, version 26.10.0].
\href{https://poppler.freedesktop.org/}{Poppler project}.

\bibitem[PyMuPDF contributors(n.d.)]{pymupdf2026}
PyMuPDF contributors. (n.d.).
\emph{PyMuPDF}, [Python package, versions 2.1 and 1.26.6].
\href{https://github.com/pymupdf/PyMuPDF}{PyMuPDF project}.

\bibitem[Python Software Foundation(n.d.)]{python2026}
Python Software Foundation. (n.d.).
\emph{Python}, [Programming language, version 3.14.10].
\href{https://www.python.org/}{Python project}.

\bibitem[PyYAML contributors(n.d.)]{pyyaml2026}
PyYAML contributors. (n.d.).
\emph{PyYAML}, [Python package, version 6.0.3].
\href{https://pyyaml.org/}{PyYAML project}.

\bibitem[R Core Team(2026)]{rcore2026}
R Core Team. (2026).
\emph{R: A language and environment for statistical computing}, [Version 4.5.3]. R Foundation for Statistical Computing, Vienna, Austria.
\href{https://www.r-project.org/}{R: A language and environment for statistical computing project}. \href{https://cran.r-project.org/src/base/R-4/}{CRAN R 4 source archive}.

\bibitem[SAS Institute(n.d.)]{sasmixed}
SAS Institute. (n.d.).
\emph{The MIXED procedure}, [SAS/STAT 15.2 User’s Guide].
\href{https://go.documentation.sas.com/api/docsets/statug/15.2/content/mixed.pdf}{SAS/STAT 15.2 User’s Guide}.

\bibitem[Satterthwaite(1946)]{satterthwaite1946}
Satterthwaite, F. E. (1946). An approximate distribution of estimates of variance components.
\emph{Biometrics Bulletin}, 2(6), 110--114.
\href{https://doi.org/10.2307/3002019}{doi:10.2307/3002019}.

\bibitem[Serlin(2000)]{serlin2000}
Serlin, R. C. (2000). Testing for robustness in Monte Carlo studies.
\emph{Psychological Methods}, 5(2), 230--240.
\href{https://doi.org/10.1037/1082-989X.5.2.230}{doi:10.1037/1082-989X.5.2.230}.

\bibitem[Tectonic contributors(n.d.)]{tectonic2026}
Tectonic contributors. (n.d.).
\emph{Tectonic}, [TeX and LaTeX engine, version 0.17.0].
\href{https://github.com/tectonic-typesetting/tectonic}{Tectonic project}.

\bibitem[Virtanen et~al.(2020)]{virtanen2020}
Virtanen, P., Gommers, R., Oliphant, T. E., et al. (2020). SciPy 1.0: Fundamental algorithms for scientific computing in Python.
\emph{Nature Methods}, 17(3), 261--272.
\href{https://doi.org/10.1038/s41592-019-0686-2}{doi:10.1038/s41592-019-0686-2}.

\end{thebibliography}
\end{document}
